\begin{figure}[htbp]
\centering
\begin{tikzpicture}[>=Stealth, every node/.style={font=\small}]
  \node (start) at (0,0) {$[\tau_0]=[\pi]$};
  \node (tail) at (2,0) {$\cdots$};
  \node (cycle) at (3.7,0) {$[\tau_l]$};
  \node (middle) at (5.4,0) {$\cdots$};
  \node (repeat) at (7.4,0) {$[\tau_{l+p}]$};
  \draw[->] (start) -- (tail);
  \draw[->] (tail) -- (cycle);
  \draw[->] (cycle) -- (middle);
  \draw[->] (middle) -- (repeat);
  \draw[dashed] (cycle.north) to[out=55,in=125]
     node[above] {same equivalence class} (repeat.north);
  \node at (1.7,-0.65) {initial segment};
  \node at (5.55,-0.65) {$p$ predecessor steps};
\end{tikzpicture}
\caption{A chain inside one finite set $D$. An arrow from $[\tau_i]$ to
$[\tau_{i+1}]$ means $\tau_i\prec\tau_{i+1}\circ\theta$.
Repetition gives $\tau_l\simeq\tau_l\circ\theta^p$ by equality of
dimensions. The proof then shows that each pullback of $\tau_l$ is
irreducible, allowing this periodicity to pass back to $\pi$.
The initial segment can be empty.}
\label{fig:finite-chain}
\end{figure}
